\section{Reflection and the quadratic large sieve}\label{sec:reflection}

We prove the second-moment estimate in \Cref{prop:reflection}.  Following
the reflection method of \cite[Part~I, Sections~5--7]{QuasiRH}, we realize
\(B_m\) as a theta average and apply the rational Weyl element.  At a
valuation-one row prime the reflected coefficient contributes a first power
of the sextic symbol; a squarefree column prime contributes power \(-4\)
in the same pairing.  Their product has exponent \(1-4\equiv3\pmod6\),
so a quadratic large sieve applies.  The reflected norm gives a column
length of order \(D^2P/Z\), and the rapid decay of the norm profile controls
the remaining terms.

\subsection{An automorphic expression for \texorpdfstring{\(B_m\)}{B m}}

All implied constants in this section may depend on the fixed field, \(S\),
the target character \(\eta\), the compact set \(\mathscr K\), and the
chosen finite family of local theta data.

We first replace elements by ideal rows.  Choose a balanced allowed generator
\(m_0\) of each exterior ideal \(\mathfrak m\).  For every nonzero
\(m\in R\) generating \(\mathfrak m\), write
\(m=\varepsilon_m m_0\) with \(\varepsilon_m\in U_S\).  When
\(m\in\mathscr B_Z(\mathscr K)\), the balanced shapes of \(m_0\), the
fixed shape set \(\mathscr K\), and
\(q_{\mathfrak m}\asymp_{\mathscr K}Z^{M_0}\) put \(\varepsilon_m\) in a
fixed compact subset of the norm-one group.  The \(S\)-unit theorem makes
\(U_S\) discrete there, so one ideal is generated by only boundedly many
elements of the row set.  Put
\[
 \Gamma=U_S^6,\qquad U_6=U_S/\Gamma,\qquad
 \upsilon_m=\varepsilon_m\Gamma\in U_6.
\]
The group \(U_6\) is finite.  The value of \(B_m\) depends on the unit
\(\varepsilon_m\) only through \(\upsilon_m\), which is fixed when the row
\(m\) is fixed.

For a nonzero row let \(j_{\mathfrak p}=v_{\mathfrak p}(\mathfrak m)\bmod6\)
in \(\{0,\ldots,5\}\).  For an allowed generator \(A\) of an integral
exterior ideal \(\mathfrak A\), reciprocity gives
\begin{equation}\label{eq:reflection-moving-twist}
 \chi_{\mathfrak A}(m)
 =\chi_{\mathfrak A}(\varepsilon_m)
   \prod_{\mathfrak p\mid\mathfrak m}
       \chi_{\mathfrak p}(A)^{j_{\mathfrak p}}.
\end{equation}
For coprime ideals this is the multiplicative reciprocity law from the
arithmetic section.  At a shared prime both sides vanish by zero extension,
including when \(j_{\mathfrak p}=0\); the zeroth power there is the unit
indicator.

Define
\(\lambda_{\varepsilon_m}(y)=H_S(\varepsilon_m,y)^{-1}\) on \(F_S^\times\).
By \Cref{eq:unit-symbol}, its value at an allowed generator \(A\) is
\(\chi_{\mathfrak A}(\varepsilon_m)\), and isotropy of \(E\) makes it
trivial on \(U_S\).  It depends only on \(\upsilon_m\).  With the character
\(L_S\) from \Cref{eq:vartheta-s}, put
\(L_{\upsilon_m}=L_S\lambda_{\varepsilon_m}\).  Then
\begin{equation}\label{eq:reflection-S-character}
 L_{\upsilon_m}(A)
   =\vartheta(\mathfrak A)\chi_{\mathfrak A}(\varepsilon_m)
 \qquad\text{for every allowed generator }A.
\end{equation}
These characters form a fixed finite family.  They are smooth, unitary,
and trivial on \(U_S\); their radial components at infinity are trivial,
and their angular component at each complex place has power \(1\) or \(-1\).

We insert the moving twists by finite upper-unipotent translations.  At an
exterior prime \(\mathfrak p\), write
\(k_{\mathfrak p}=\mathcal O_F/\mathfrak p\) and use its allowed generator
\(p\) as a local uniformizer.  For \(1\le j\le5\) put
\[
 \tau_{\mathfrak p,j}
 =\sum_{v\in k_{\mathfrak p}^{\times}}
       \chi_{\mathfrak p}(v)^{-j}\psi_{\mathfrak p}(v/p).
\]
The character is nontrivial, so
\(\lvert\tau_{\mathfrak p,j}\rvert=q_{\mathfrak p}^{1/2}\).  Define
\begin{equation}\label{eq:reflection-Fourier-mask}
 \alpha_{\mathfrak p,j}(v)=
 \begin{cases}
  \tau_{\mathfrak p,j}^{-1}\chi_{\mathfrak p}(v)^{-j},
       &1\le j\le5,\ v\ne0,\\
  0,   &1\le j\le5,\ v=0,\\
  1-q_{\mathfrak p}^{-1},&j=0,\ v=0,\\
  -q_{\mathfrak p}^{-1},&j=0,\ v\ne0.
 \end{cases}
\end{equation}
For every \(t\in\mathcal O_{\mathfrak p}\),
\begin{equation}\label{eq:reflection-mask-identity}
 \sum_{v\in k_{\mathfrak p}}\alpha_{\mathfrak p,j}(v)
                \psi_{\mathfrak p}(tv/p)
       =\chi_{\mathfrak p}(t)^j.
\end{equation}
For \(j\ne0\), substitution by \(t^{-1}\) proves this for a unit \(t\),
and both sides vanish for a nonunit.  For \(j=0\), the left side is
\(1-q_{\mathfrak p}^{-1}\sum_v\psi_{\mathfrak p}(tv/p)\), the unit indicator.
Thus the formula includes the masked zeroth power.

For each \(\upsilon\in U_6\), choose the theta vector in
\Cref{prop:theta-input} whose finite \(S\)-Whittaker function is the
indicator of the common small product subgroup contained in sixth powers
and in \(\ker L_\upsilon\), and whose complex angular type is that of
\(L_\upsilon\).  Denote its automorphic function by \(\theta_\upsilon\).
The cutoff and the normalizations are those of
\Cref{eq:theta-chosen-S-function}; these choices form a fixed finite family.
For the remainder of this subsection fix a nonzero \(m\), put
\(\upsilon=\upsilon_m\), and write \(L=L_\upsilon\).

For \(y\in F_S^\times\) and digits
\(\boldsymbol v=(v_{\mathfrak p})_{\mathfrak p\mid\mathfrak m}\), let
\(g(y,\boldsymbol v)\) have component \(a(y)\) at \(S\), component
\(n(v_{\mathfrak p}/p)\) at \(\mathfrak p\mid\mathfrak m\), and identity
at the other exterior places.  Use the lifts in
\Cref{eq:theta-lift-identities}, choose representatives of the residue
digits, and set
\[
 \boldsymbol\alpha_m(\boldsymbol v)
   =\prod_{\mathfrak p\mid\mathfrak m}
          \alpha_{\mathfrak p,j_{\mathfrak p}}(v_{\mathfrak p}),
 \qquad w_+=1,\quad w_-=w.
\]
The combined digit sums that we shall integrate are
\begin{equation}\label{eq:reflection-combined-digits}
\begin{split}
 \mathscr D_m^\pm(y)
   &=\sum_{\boldsymbol v}\boldsymbol\alpha_m(\boldsymbol v)
          \theta_\upsilon(w_\pm g(y,\boldsymbol v)),\\
 \mathscr D_{m,N}^\pm(y)
   &=\sum_{\boldsymbol v}\boldsymbol\alpha_m(\boldsymbol v)
          \theta_{\upsilon,N}(w_\pm g(y,\boldsymbol v)).
\end{split}
\end{equation}
Rational automorphy gives \(\mathscr D_m^+=\mathscr D_m^-\).  We next
show that all four sums descend under \(\Gamma\).

For \(\gamma\in\Gamma\), choose representatives
\(v'_{\mathfrak p}\equiv\gamma v_{\mathfrak p}\pmod{\mathfrak p}\), and put
\[
 u_{\mathfrak p}
   =(\gamma v_{\mathfrak p}-v'_{\mathfrak p})/p\in\mathcal O_{\mathfrak p}.
\]
Let \(k_S=1\),
\(k_{\mathfrak p}=n(u_{\mathfrak p})a(\gamma)\) for
\(\mathfrak p\mid\mathfrak m\), and \(k_{\mathfrak q}=a(\gamma)\) at the
other exterior places.  Because \(\gamma\) is an exterior unit, these are
compact lifts fixing the original exterior spherical vector.  At \(S\)
the torus factors involving \(\gamma\) are trivial, since it is a sixth
power.  At a row prime the exact root relation gives
\[
 n(v'_{\mathfrak p}/p)n(u_{\mathfrak p})a(\gamma)
       =a(\gamma)n(v_{\mathfrak p}/p).
\]
The rational diagonal and Weyl lifts therefore satisfy
\begin{equation}\label{eq:reflection-quotient-lifts}
\begin{split}
 g(\gamma y,\boldsymbol v')
   &=a(\gamma)g(y,\boldsymbol v)k^{-1},\\
 w g(\gamma y,\boldsymbol v')
   &=a(\gamma^{-1})w g(y,\boldsymbol v)k^{-1}.
\end{split}
\end{equation}
The second identity has the same compact factor at the far right.

The upper constant term has the required invariance at these right
arguments as well.  For \(q\in F^\times\), any adelic \(h\), and any such
exterior compact \(k\), the root relation and the product formula give
\begin{equation}\label{eq:reflection-constant-term-descent}
\begin{split}
 \theta_{\upsilon,N}(a(q)hk^{-1})
 &=\int_{F\backslash\Adele_F}
       \theta_\upsilon(a(q)n(x/q)hk^{-1})\,\dd x\\
 &=\theta_{\upsilon,N}(h).
\end{split}
\end{equation}
Indeed \(n(x)a(q)=a(q)n(x/q)\); rational automorphy removes \(a(q)\),
and right invariance removes \(k^{-1}\) at the far right.  The substitution
\(x\mapsto x/q\) preserves Haar measure on \(F\backslash\Adele_F\).
This argument applies in particular when \(h_{\mathfrak p}=wn(v/p)\).
Finally \(\chi_{\mathfrak p}(\gamma)=1\), so
\(\alpha_{\mathfrak p,j}(\gamma v)=\alpha_{\mathfrak p,j}(v)\), including
the zero digit.  Reindexing by \(\boldsymbol v'=\gamma\boldsymbol v\) in
\Cref{eq:reflection-quotient-lifts,eq:reflection-constant-term-descent}
proves the asserted descent of the combined sums before and after \(w\).

Recall \(X_l=\{y\in F_S^\times:|y|_S=l\}\) and the sheet measures from
\Cref{eq:sheet-disintegration}.  The quotient \(X_l/\Gamma\) is compact;
choose a measurable fundamental set \(\Omega_1\subset X_1\) and put
\(\Omega_l=\ell_l\Omega_1\), using the norm section from that disintegration.
Every integral over \(X_l/\Gamma\) means the integral over \(\Omega_l\)
with the restricted sheet measure, equivalently the quotient measure
induced using counting measure on \(\Gamma\).  The factors
\(\one_B([y])\) and \(L(y)^{-1}\) also descend, since
\([\gamma]=1\) in \(\mathcal V\) and \(L(\gamma)=1\).  The resulting
integrands are integrable on the quotient: the digit sums are finite and
the factors are smooth or locally constant.

Both combined constant terms have zero weighted integral.  To see this,
rotate \(y\) by a complex number \(z\) of modulus one at one complex place.
In the plus sum this
acts by the left diagonal \(a(z)\); in the minus sum it acts by \(a(z^{-1})\).
The constant term belongs to the torus induction in
\Cref{prop:theta-input}.  Its complex inducing law is radial for every
right argument, so both sums are unchanged, including at the exterior
arguments \(wn(v/p)\).  The class condition is unchanged because
\(\C^\times\) is sixth-divisible, while \(L(y)^{-1}\) transforms by a
nontrivial angular character.  Haar invariance therefore gives
\begin{equation}\label{eq:reflection-angular-cancellation}
 \int_{X_l/\Gamma}\one_B([y])L(y)^{-1}
             \mathscr D_{m,N}^\pm(y)\,\dd_l^\times y=0.
\end{equation}
We may now define the quotient average
\begin{equation}\label{eq:reflection-theta-average}
\begin{split}
 \mathcal T_m(l)
  &=\int_{X_l/\Gamma}\one_B([y])L(y)^{-1}
             \mathscr D_m^+(y)\,\dd_l^\times y\\
  &=\int_{X_l/\Gamma}\one_B([y])L(y)^{-1}
             \mathscr D_m^-(y)\,\dd_l^\times y.
\end{split}
\end{equation}

\begin{lemma}[Theta representation of the smoothed sum]
\label{lem:reflection-theta-representation}
Let \(V\) and \(\mathcal K\) be the common unrotated profile and Gaussian
averaging weight from \Cref{lem:gaussian-smoothing}.  For \(l,Z>0\),
\begin{align}
 \mathcal T_m(l)
 &=\sum_{\substack{\mathfrak c\ \mathrm{squarefree}\\\mathfrak n}}
    \frac{\gamma_2(\mathfrak c)}{q_{\mathfrak c}^{1/2}q_{\mathfrak n}}
    L(A)\prod_{\mathfrak p\mid\mathfrak m}
                  \chi_{\mathfrak p}(A)^{j_{\mathfrak p}}
    V(q_{\mathfrak A}l),                                      \label{eq:reflection-unreflected}\\
 B_m(Z)
 &=\int_0^\infty\mathcal K(T)\mathcal T_m(T/Z)\,\frac{\dd T}{T}.
                                                               \label{eq:reflection-exact-average}
\end{align}
Here \(\mathfrak A=\mathfrak c\mathfrak n^3\) and \(A=cn^3\) is an
allowed product generator.  The ideal sums are over integral exterior
ideals.  Both formulas converge absolutely after the sheet integrations,
and the second converges absolutely after the \(T\)-integration.
\end{lemma}

\begin{proof}
Use the Fourier expansion and normalized factorization in
\Cref{eq:theta-fourier-expansion,eq:theta-normalized-factorization} in the
plus average.  The constant term vanishes by
\Cref{eq:reflection-angular-cancellation}.  At a nonconstant rational
frequency \(\nu\), a row-prime factor is
\[
 W_{\mathfrak p}(a(\nu)n(v/p))
   =\psi_{\mathfrak p}(\nu v/p)W_{\mathfrak p}(a(\nu)).
\]
Spherical support forces \(\nu\in R\); summing the digits with
\Cref{eq:reflection-mask-identity} supplies
\(\prod_{\mathfrak p\mid\mathfrak m}\chi_{\mathfrak p}(\nu)^{j_{\mathfrak p}}\).

Write \(f_\upsilon(v)\) for the chosen unrotated \(S\)-Whittaker function.
The \(S\)-factor of the frequency is
\(\mathcal C_S(\nu,y)f_\upsilon(\nu y)\).  Its cutoff has
\([\nu y]=1\) in \(\mathcal V\).  Together with \([y]\in B\), this gives
\([\nu]\in B\).  By alternation and \Cref{eq:reflection-S-cocycle},
\[
 \mathcal C_S(\nu,y)
   =H_S(\nu,\nu^{-1})^2H_S(\nu,\nu y)^2=1
\]
on that support.  If \(A\) is an allowed generator of the exterior ideal
of \(\nu\), then \(\nu/A\in U_S\).  The direct sum \(E\oplus B\) and the
injection \(U_6\longrightarrow E\) now imply \(\nu=\gamma A\) for
\(\gamma\in\Gamma\).  The grouped exterior factor is unchanged by this
multiplication, by the local unit rule and
\(\chi_{\mathfrak p}(\gamma)=1\).

We justify the orbit substitution with an absolute integral first.
The sets \(\gamma\Omega_l\) partition \(X_l\) up to null sets.  Tonelli,
multiplicative Haar invariance, and \Cref{eq:sheet-weighted-bound} give
for \(\sigma>1/2\)
\begin{equation}\label{eq:reflection-forward-absolute}
\begin{split}
 \sum_{\gamma\in\Gamma}\int_{\Omega_l}
       |f_\upsilon(\gamma Ay)|\,\dd_l^\times y
  &=\int_{X_{q_{\mathfrak A}l}}|f_\upsilon(v)|\,
                      \dd_{q_{\mathfrak A}l}^\times v\\
  &\ll_\sigma(q_{\mathfrak A}l)^{-\sigma}.
\end{split}
\end{equation}
The signed substitution \(v=\gamma Ay\) is therefore legitimate.  Since
\([y]=[A]^{-1}[v]\) and \(L(y)^{-1}=L(A)L(v)^{-1}\), it gives
\begin{equation}\label{eq:reflection-forward-unfolding}
\begin{split}
 &\sum_{\gamma\in\Gamma}\int_{\Omega_l}
       \one_B([y])L(y)^{-1}f_\upsilon(\gamma Ay)\,\dd_l^\times y\\
 &\quad=L(A)\int_{X_{q_{\mathfrak A}l}}
       \one_B([A]^{-1}[v])L(v)^{-1}f_\upsilon(v)\,
                    \dd_{q_{\mathfrak A}l}^\times v\\
 &\quad=L(A)V(q_{\mathfrak A}l).
\end{split}
\end{equation}
For the last equality, \([v]=1\) on the cutoff, \([A]\in B\), and the
angular characters cancel as in \Cref{eq:theta-V-profile}.  Thus the
integral depends on \(A\) only through its norm.  Multiplicative Haar
measure introduces no norm Jacobian.

By \Cref{cor:theta-exterior-coefficient}, the exterior spherical
coefficient vanishes unless \(\mathfrak A=\mathfrak c\mathfrak n^3\)
with \(\mathfrak c\) squarefree, and then equals
\(\gamma_2(\mathfrak c)/(q_{\mathfrak c}^{1/2}q_{\mathfrak n})\).
This includes primes shared by \(\mathfrak c\) and \(\mathfrak n\).
Together with \Cref{eq:reflection-forward-unfolding} it proves
\Cref{eq:reflection-unreflected}.

Here are the absolute bounds for the exchanges.  Summing
\Cref{eq:reflection-forward-absolute} with the exterior magnitudes is at
most \(l^{-\sigma}\) times the convergent expression in
\Cref{eq:reflection-original-dirichlet-majorant}.  Before evaluating
the finite digit sums, their total absolute mass is at most
\(q_{\mathfrak p}^{1/2}\) for \(j\ne0\) and is less than \(2\) for \(j=0\).
For a fixed row their finite product multiplies the same spherical
majorant, which justifies that earlier exchange.  Finally
\Cref{lem:gaussian-smoothing} gives every absolute moment of \(\mathcal K\).
The averaged absolute sum is bounded by a constant times
\[
 Z^\sigma M_{-\sigma}(\mathcal K)
\]
times the finite sum in \Cref{eq:reflection-original-dirichlet-majorant}.
We may therefore apply \Cref{eq:theta-exact-gaussian} term by term.
\Cref{eq:reflection-moving-twist,eq:reflection-S-character} and the
definition of \(B_m\) give \Cref{eq:reflection-exact-average}.
\end{proof}

\subsection{The local Weyl calculation}

The second expression in \Cref{eq:reflection-theta-average} is now the
starting point.  At \(S\), \(wa(y)=a(y^{-1})w\) replaces the vector by a
fixed rotation and the argument by \(y^{-1}\).  Its constant term has
already been removed in \Cref{eq:reflection-angular-cancellation}.
At an exterior place with zero or no digit, \(w\) is compact.  We compute
the effect of each nonzero digit.

\begin{lemma}[Local effect of reflection]\label{lem:reflection-local-Weyl}
Let \(\mathfrak p\) be an exterior prime and put
\(W_e=W_{\mathfrak p}(a(p^e))\).  For \(j\in\{0,\ldots,5\}\), a unit \(t\),
and an integer \(e\), define
\[
 \mathcal R_{\mathfrak p,j,e}(t)
 =\sum_{v\in k_{\mathfrak p}^{\times}}\alpha_{\mathfrak p,j}(v)
       W_{\mathfrak p}(a(p^{e-2}t)w n(v/p)).
\]
It vanishes for \(e<0\).  For \(e\ge0\),
\begin{equation}\label{eq:reflection-local-formula}
 \mathcal R_{\mathfrak p,j,e}(t)
 =W_e\chi_{\mathfrak p}(t)^{4+2e}
   \sum_{v\in k_{\mathfrak p}^{\times}}
     \alpha_{\mathfrak p,j}(v)\chi_{\mathfrak p}(v)^4
     \psi_{\mathfrak p}(p^{e-1}t/v).
\end{equation}
For \(j=1\), it vanishes for \(e>0\), while
\begin{equation}\label{eq:reflection-valuation-one}
\begin{split}
 \mathcal R_{\mathfrak p,1,0}(t)
      &=\omega_{\mathfrak p}\chi_{\mathfrak p}(t),
       \qquad |\omega_{\mathfrak p}|=1,\\
 \omega_{\mathfrak p}
 &=\frac{\sum_{x\ne0}\chi_{\mathfrak p}(x)^3\psi_{\mathfrak p}(x/p)}
        {\tau_{\mathfrak p,1}}.
\end{split}
\end{equation}
For all \(j\), its magnitude at \(e=0\) is at most \(1\), and for \(e>0\)
is at most \(q_{\mathfrak p}^{1/2}|W_e|\).  More precisely, for \(e\ge0\)
put \(s_{\mathfrak p,j,e}=\mathcal R_{\mathfrak p,j,e}(1)\) and choose
the exponents modulo \(6\) by
\begin{equation}\label{eq:reflection-local-monomial}
 b_{j,0}\equiv8-j,\qquad b_{j,e}\equiv4+2e\ (e>0),\qquad
 \mathcal R_{\mathfrak p,j,e}(t)
       =s_{\mathfrak p,j,e}\chi_{\mathfrak p}(t)^{b_{j,e}}.
\end{equation}
For \(e>0\), \(s_{\mathfrak p,j,e}=0\) unless \(j=4\), and it also
vanishes for \(e\equiv2\pmod3\).  The zero-digit term is available only
for \(j=0\); at a frequency \(\nu\) it is
\((1-q_{\mathfrak p}^{-1})W_{\mathfrak p}(a(\nu))\), with ordinary
spherical support.
\end{lemma}

\begin{proof}
For a unit digit \(v\), the exact lift identity in
\Cref{eq:theta-lift-identities}, with \(H=v/p\), is
\[
 w n(H)=n(H^{-1})a(H^{-2})k_v,\qquad
 k_v=a(-1)\bar n(p/v)
\]
with \(k_v\) in the spherical compact group.  Only the diagonal minus
sign has been absorbed in \(k_v\); its symbols are trivial because
\(F\) contains \(\mu_{12}\).  Commuting the upper root past
\(a(p^{e-2}t)\) gives the positive additive argument
\(\psi_{\mathfrak p}(p^{e-1}t/v)\).  The effective torus law gives, for
integers \(e_1,e_2\) and units \(u,v\),
\[
 \mathcal C_{\mathfrak p}(p^{e_1}u,p^{e_2}v)
   =\chi_{\mathfrak p}(u)^{2e_2}\chi_{\mathfrak p}(v)^{-2e_1}.
\]
The unit--unit and uniformizer--uniformizer factors are trivial.  Thus
torus multiplication contributes
\[
 \mathcal C_{\mathfrak p}(p^{e-2}t,p^2v^{-2})
       =\chi_{\mathfrak p}(t)^4\chi_{\mathfrak p}(v)^{4(e-2)}.
\]
The remaining spherical argument is \(p^etv^{-2}\).  It is zero for
\(e<0\); otherwise \Cref{eq:theta-local-values} contributes
\(\chi_{\mathfrak p}(t)^{2e}\chi_{\mathfrak p}(v)^{-4e}W_e\).
Their product proves \Cref{eq:reflection-local-formula}.

For \(e>0\) the additive character is trivial.  Unit-character
orthogonality makes the digit sum zero unless \(j=4\), including the
case \(j=0\).  For every \(j\) its absolute value is at most
\(q_{\mathfrak p}^{1/2}\): use
\(|\alpha_{\mathfrak p,j}(v)|=q_{\mathfrak p}^{-1/2}\) for \(j\ne0\)
and \(q_{\mathfrak p}^{-1}\) for \(j=0\).  Its unit dependence is the
power \(4+2e\), and \(W_e=0\) for \(e\equiv2\pmod3\).

For \(e=0\), substitute \(x=t/v\).  If \(j\ne0\), this leaves the power
\(\chi_{\mathfrak p}(t)^{8-j}\), divided by \(\tau_{\mathfrak p,j}\),
and the Gauss sum of \(\chi_{\mathfrak p}^{j-4}\) against
\(\psi_{\mathfrak p}(x/p)\).  That sum has magnitude
\(q_{\mathfrak p}^{1/2}\), except for \(j=4\), when it has magnitude \(1\).
For \(j=1\) its character is quadratic and \(8-j\equiv1\pmod6\),
giving \Cref{eq:reflection-valuation-one}.  For \(j=0\), the same
substitution leaves \(\chi_{\mathfrak p}(t)^8\) times a nontrivial
Gauss sum divided by \(-q_{\mathfrak p}\), again of magnitude at most one.
These observations prove \Cref{eq:reflection-local-monomial} and the
remaining bounds.  Finally, compactness of \(w\) gives the zero-digit
assertion by spherical right invariance.
\end{proof}

For a row \(\mathfrak m=\mathfrak h\mathfrak d\), call a prime active
when its digit is nonzero.  Every prime of \(\mathfrak d\) is active,
because there \(j_{\mathfrak p}=1\).  For \(\mathfrak p\mid\mathfrak h\)
put \(j_{\mathfrak p}=v_{\mathfrak p}(\mathfrak h)\bmod6\).  The exact
set of possible active products in \(\mathfrak h\) is
\begin{equation}\label{eq:reflection-active-set}
 \mathcal F(\mathfrak h)=
 \left\{\mathfrak f\ \mathrm{squarefree}:
   \mathfrak f\mid\rad(\mathfrak h),\
   \mathfrak p\mid\mathfrak f\ \text{if }\
       \mathfrak p\mid\mathfrak h,\ j_{\mathfrak p}\ne0\right\}.
\end{equation}
Only a prime with \(j_{\mathfrak p}=0\) can be inactive.  For
\(\mathfrak f\in\mathcal F(\mathfrak h)\), its exact inactive scalar is
\[
 I_{\mathfrak h,\mathfrak f}
    =\prod_{\substack{\mathfrak p\mid\mathfrak h\\
                      \mathfrak p\nmid\mathfrak f}}
           (1-q_{\mathfrak p}^{-1}),
\]
and powerfulness gives
\begin{equation}\label{eq:reflection-active-norm}
 \mathfrak f\mid\rad(\mathfrak h),\qquad
 q_{\mathfrak f}^2\le q_{\mathfrak h}.
\end{equation}
Let \(\mathcal A(\mathfrak f)\) be the integral exterior ideals
\(\mathfrak a'_{\mathfrak f}\) supported on \(\mathfrak f\), with every
exponent congruent to \(0\) or \(1\) modulo \(3\); set
\(\mathcal A(1)=\{1\}\).  This is a convenient common support: the actual
coefficient \(s_{\mathfrak p,j_{\mathfrak p},e}\) remains zero in the
locally vanishing cases of \Cref{lem:reflection-local-Weyl}.

The local support shows that every frequency with a nonzero grouped
exterior coefficient occurs among
\begin{equation}\label{eq:reflection-frequency}
 \nu=\varepsilon\frac{a'_{\mathfrak f}c'(n')^3}{(df)^2},
 \qquad
 \mathfrak a'=\mathfrak a'_{\mathfrak f}\mathfrak c'(\mathfrak n')^3,
 \qquad \varepsilon\in U_S,
\end{equation}
where \(\mathfrak a'_{\mathfrak f}\in\mathcal A(\mathfrak f)\),
\(\mathfrak c'\) is squarefree, and
\((\mathfrak c'\mathfrak n',\mathfrak d\mathfrak f)=1\).
The displayed elements \(d,f,a'_{\mathfrak f},c',n'\) are allowed
generators.  The ideals \(\mathfrak c'\) and \(\mathfrak n'\) may share
primes, and primes of \(\mathfrak h\) omitted from \(\mathfrak f\) may
occur in either.  The valuation-one formula forces numerator exponent
zero at each prime of \(\mathfrak d\).  Away from
\(\mathfrak d\mathfrak f\), the valuations \(3k\) and \(3k+1\)
determine \(\mathfrak c'\) and \(\mathfrak n'\) uniquely.  The remaining
ambiguity is the displayed \(S\)-unit, since allowed product generators
are unique modulo \(\Gamma\).

For later absolute estimates define
\begin{equation}\label{eq:reflection-frozen-weight}
 \beta_{\mathfrak p}(e)=
 \begin{cases}
  1,&e=0,\\
  q_{\mathfrak p}^{-k},&e=1+3k,\ k\ge0,\\
  q_{\mathfrak p}^{1/2-k},&e=3k,\ k\ge1,\\
  0,&\text{otherwise},
 \end{cases}
 \qquad
 \beta_{\mathfrak f}(\mathfrak a'_{\mathfrak f})
    =\prod_{\mathfrak p\mid\mathfrak f}
       \beta_{\mathfrak p}(v_{\mathfrak p}(\mathfrak a'_{\mathfrak f})).
\end{equation}
The local lemma and the spherical magnitudes give
\(|s_{\mathfrak p,j,e}|\le\beta_{\mathfrak p}(e)\).  The other exterior
magnitudes together are at most
\(q_{\mathfrak c'}^{-1/2}q_{\mathfrak n'}^{-1}\).

\subsection{The full reflected expansion}

We now assemble these local factors.  The \(S\)-integrals form a fixed
family of one-variable profiles.  Let
\(f_{w,\upsilon}(v)\) be the \(S\)-Whittaker function of the chosen vector
after the fixed Weyl rotation.  Choose once a presentation
\[
 f_{w,\upsilon}(v)
   =\sum_{\ell=1}^{J_\upsilon}\prod_{a\in S}
          f_{\upsilon,\ell,a}(v_a).
\]
The finite \(S\)-functional and the fixed cutoff vector provide such
presentations by \Cref{prop:theta-input,eq:theta-chosen-S-function};
their lengths are bounded over
the finite row-class family.  All fixed scalar coefficients are included
in the displayed local factors.  Write \(H_S\) also for its induced
pairing on \(\mathcal V\).  For \(\kappa\in\mathcal V\), define
\begin{equation}\label{eq:reflection-profiles}
\begin{split}
 H_{\upsilon,\kappa,\ell}(x)
  &=\int_{X_x}\one_B(\kappa[v]^{-1})H_S(\kappa,[v])^2L_\upsilon(v)\\
  &\hspace{25mm}\times
       \prod_{a\in S}f_{\upsilon,\ell,a}(v_a)\,\dd_x^\times v,\\
 \Phi_{\upsilon,\kappa,\ell}(x)
  &=\int_0^\infty\mathcal K(T)H_{\upsilon,\kappa,\ell}(x/T)
           \,\frac{\dd T}{T},\\
 \Phi_{\upsilon,\kappa}(x)
  &=\sum_{\ell=1}^{J_\upsilon}\Phi_{\upsilon,\kappa,\ell}(x).
\end{split}
\end{equation}
The class multiplier is bounded and depends only on \(\mathcal V\), by
\Cref{eq:reflection-S-cocycle}.  The local bounds for the fixed rotated
vectors and \Cref{lem:sheet-profile,lem:gaussian-smoothing} therefore give
one \(\alpha>0\) such that, for every \(N>0\) and integer \(j\ge0\),
\begin{equation}\label{eq:reflection-profile-bound}
 |(x\partial_x)^j\Phi_{\upsilon,\kappa}(x)|
       \ll_{N,j}\min(x^\alpha,x^{-N})\qquad(x>0).
\end{equation}
The constants are uniform in \(\upsilon\in U_6\) and
\(\kappa\in\mathcal V\).  This is a fixed finite family: the functional
and tensor terms have been combined in \(\Phi_{\upsilon,\kappa}\), while
the bounded class function stays inside its integral.  In particular
those fixed coefficients will not enter the local majorant
\(\beta_{\mathfrak f}\).

Choose representatives \(\varepsilon_u\) for the separate frequency-unit
classes \(u\in U_6\).  For
\[
 \mathfrak t=(\mathfrak h,\mathfrak f,\mathfrak a'_{\mathfrak f},
              \mathfrak n',u),\qquad
 \mathfrak f\in\mathcal F(\mathfrak h),\quad
 \mathfrak a'_{\mathfrak f}\in\mathcal A(\mathfrak f),\quad
 (\mathfrak n',\mathfrak f)=1,
\]
and class slices \(s_d,s_c\in B\), put
\begin{equation}\label{eq:reflection-profile-class}
 \kappa_6(\mathfrak t;s_d,s_c)
   =[\varepsilon_u][a'_{\mathfrak f}]s_c[n']^3
        s_d^{-2}[f]^{-2}\in\mathcal V.
\end{equation}
This is a selected class, rather than a further summation index.  For
\(\mathfrak p\mid\mathfrak f\), abbreviate
\[
 e_{\mathfrak p}=v_{\mathfrak p}(\mathfrak a'_{\mathfrak f}),\qquad
 s_{\mathfrak p}=s_{\mathfrak p,j_{\mathfrak p},e_{\mathfrak p}},\qquad
 b_{\mathfrak p}=b_{j_{\mathfrak p},e_{\mathfrak p}}.
\]
With \(\upsilon,\mathfrak t,s_d,s_c\) fixed, the next three factors organize the
dependence on the two varying ideals: \(\omega\) depends on neither,
\(\rho\) depends on the row \(\mathfrak d\) alone, and \(\lambda\) depends
on the column \(\mathfrak c'\) alone.  All three are independent of \(Z\).
The quadratic symbol and norm profile will contain the remaining joint
dependence on \(\mathfrak d\) and \(\mathfrak c'\).
For \(L=L_\upsilon\) and \(\varepsilon=\varepsilon_u\), define
\begin{equation}\label{eq:reflection-frozen-factor}
\begin{split}
 \omega_{\upsilon,\mathfrak t}
  ={}&I_{\mathfrak h,\mathfrak f}
       L(a'_{\mathfrak f})^{-1}L(n')^{-3}L(f)^2\\
    &{}\times\prod_{\mathfrak p\mid\mathfrak f}
       s_{\mathfrak p}
       \chi_{\mathfrak p}\!\left(
         \varepsilon\frac{a'_{\mathfrak f}}{p^{e_{\mathfrak p}}}
         (n')^3(f/p)^{-2}\right)^{b_{\mathfrak p}}.
\end{split}
\end{equation}
For a squarefree \(\mathfrak d\) with
\((\mathfrak d,\mathfrak h\mathfrak n')=1\) and \([d]=s_d\), put
\begin{equation}\label{eq:reflection-row-factor}
\begin{split}
 \rho_{\upsilon,\mathfrak t,s_d}(\mathfrak d)
  ={}&L(d)^2\\
    &{}\times\prod_{\mathfrak p\mid\mathfrak d}
        \omega_{\mathfrak p}
        \chi_{\mathfrak p}\!\left(
           \varepsilon a'_{\mathfrak f}(n')^3(d/p)^{-2}f^{-2}\right)\\
    &{}\times\prod_{\mathfrak p\mid\mathfrak f}
              \chi_{\mathfrak p}(d)^{-2b_{\mathfrak p}},
\end{split}
\end{equation}
and put it equal to zero otherwise.  For a squarefree \(\mathfrak c'\)
with \((\mathfrak c',\mathfrak f)=1\) and \([c']=s_c\), put
\begin{equation}\label{eq:reflection-column-factor}
\begin{split}
 \lambda_{\upsilon,\mathfrak t,s_c}(\mathfrak c')
  ={}&L(c')^{-1}
       \prod_{\mathfrak p\mid\mathfrak f}
                 \chi_{\mathfrak p}(c')^{b_{\mathfrak p}}\\
    &{}\times\prod_{\mathfrak q\mid\mathfrak c'}
       g_{2,\mathfrak q}
       \chi_{\mathfrak q}\!\left(
          \varepsilon a'_{\mathfrak f}(c'/q)f^{-2}\right)^2,
\end{split}
\end{equation}
and put it equal to zero otherwise.  All character arguments in these
formulas are units at their primes on the stated supports.  Thus negative
character powers are well-defined.  These definitions give one coherent
row function and one coherent column function for each frozen tuple and
class slice.

\begin{lemma}[The reflected expansion]\label{lem:reflection-quadratic-pairing}
For every nonzero \(m\in R\) and \(Z>0\), let
\((m)_S=\mathfrak h\mathfrak d\) and \(\upsilon=\upsilon_m\).  Then
\begin{equation}\label{eq:reflection-global-expansion}
\begin{split}
 B_m(Z)
  ={}&\sum_{\mathfrak f\in\mathcal F(\mathfrak h)}
      \sum_{\substack{\mathfrak a'_{\mathfrak f}\in\mathcal A(\mathfrak f)\\
                      \mathfrak n'\\(\mathfrak n',\mathfrak f)=1}}
      \sum_{\substack{u\in U_6\\s_d,s_c\in B}}
       \frac{\omega_{\upsilon,\mathfrak t}}{q_{\mathfrak n'}}\\
   &{}\times\rho_{\upsilon,\mathfrak t,s_d}(\mathfrak d)
      \Biggl[\sum_{\mathfrak c'\ \mathrm{squarefree}}
       \frac{\lambda_{\upsilon,\mathfrak t,s_c}(\mathfrak c')}
            {q_{\mathfrak c'}^{1/2}}
       \chi_{\mathfrak d}(c')^3\\
   &{}\times\Phi_{\upsilon,\kappa_6(\mathfrak t;s_d,s_c)}
       \!\left(\frac{Zq_{\mathfrak a'_{\mathfrak f}}
                     q_{\mathfrak c'}q_{\mathfrak n'}^3}
                    {q_{\mathfrak d}^2q_{\mathfrak f}^2}\right)\Biggr].
\end{split}
\end{equation}
All ideal sums are over integral exterior ideals, and the expansion is
absolutely convergent.  The coefficients satisfy
\[
 |\omega_{\upsilon,\mathfrak t}|
     \le\beta_{\mathfrak f}(\mathfrak a'_{\mathfrak f}),\qquad
 |\rho_{\upsilon,\mathfrak t,s_d}|\le1,\qquad
 |\lambda_{\upsilon,\mathfrak t,s_c}|\le1.
\]
The row-unit class \(\upsilon\) is fixed;
\(u\) is the independently summed frequency-unit class.
\end{lemma}

\begin{proof}
We use the minus average in \Cref{eq:reflection-theta-average} and its
nonconstant Fourier terms.  We first supply the absolute bound needed
before unfolding any signed unit orbit.  Define
\[
 F^{\mathrm{abs}}_{w,\upsilon}(v)
    =\sum_{\ell=1}^{J_\upsilon}
           \prod_{a\in S}|f_{\upsilon,\ell,a}(v_a)|,\qquad
 H^{\mathrm{abs}}_{w,\upsilon}(x)
    =\int_{X_x}F^{\mathrm{abs}}_{w,\upsilon}(v)\,\dd_x^\times v.
\]
This is an absolute profile of \Cref{lem:sheet-profile}; in particular
\(H^{\mathrm{abs}}_{w,\upsilon}(x)\ll_\sigma x^{-\sigma}\) for
\(\sigma>1/2\), uniformly over the fixed row-class family.  For any
\(\nu_0\in F^\times\), the sets \(\gamma\Omega_l^{-1}\) partition
\(X_{1/l}\).  Tonelli and multiplicative Haar invariance give
\begin{equation}\label{eq:reflection-reflected-absolute}
\begin{split}
 \sum_{\gamma\in\Gamma}\int_{\Omega_l}
       F^{\mathrm{abs}}_{w,\upsilon}(\gamma\nu_0/y)\,\dd_l^\times y
  &=H^{\mathrm{abs}}_{w,\upsilon}(|\nu_0|_S/l)\\
  &\ll_\sigma(|\nu_0|_S/l)^{-\sigma}.
\end{split}
\end{equation}

For a fixed row, this also justifies the exchanges before evaluating the
digits.  At one nonzero digit, the absolute Whittaker factor in the proof of
\Cref{lem:reflection-local-Weyl} is \(|W_e|\) at numerator valuation
\(e\ge0\).  Summing the absolute digit weights costs at most
\(q_{\mathfrak p}^{1/2}\) for \(j\ne0\) and is less than \(2\) for \(j=0\).
For \(\sigma>1/2\), the weighted spherical series is
\[
 \sum_{e\ge0}|W_e|q_{\mathfrak p}^{-\sigma e}
   =\frac{1+q_{\mathfrak p}^{-1/2-\sigma}}
          {1-q_{\mathfrak p}^{-1-3\sigma}}<\infty.
\]
The denominator shift \(e-2\) multiplies this by
\(q_{\mathfrak p}^{2\sigma}\).  There are only finitely many row primes,
so their resulting product is finite for the fixed row.  Here the numerator
valuations at the primes of \(\mathfrak d\) are still summed.  The product
of ordinary spherical series at the other primes is bounded by
\(\zeta_F^S(\sigma+1/2)\zeta_F^S(3\sigma+1)\).  Partitioning the units
into the finitely many \(U_6\)-classes and applying
\Cref{eq:reflection-reflected-absolute} bounds the absolute sheet sum.
The \(T\)-average then uses \(M_\sigma(\mathcal K)<\infty\).
Thus these preliminary exchanges are valid for each fixed row.  The
uniform estimates below use the coefficients after the finite digit sums
have been evaluated.

Partition the zero and nonzero digits by
\(\mathfrak f\in\mathcal F(\mathfrak h)\), apply
\Cref{lem:reflection-local-Weyl}, and parameterize the frequencies by
\Cref{eq:reflection-frequency}.  For the fixed representative
\(\varepsilon=\varepsilon_u\), write
\[
 \nu=\gamma\nu_0,\qquad
 \nu_0=\varepsilon\frac{a'_{\mathfrak f}c'(n')^3}{(df)^2},
 \qquad \gamma\in\Gamma.
\]
The valuations and then the frequency-unit class give a unique such
parameterization on the stated support.  At
\(\mathfrak p\mid\mathfrak d\), the local unit of \(\nu_0=p^{-2}t\) is
\[
 t=\varepsilon a'_{\mathfrak f}c'(n')^3(d/p)^{-2}f^{-2}.
\]
At \(\mathfrak p\mid\mathfrak f\), the local unit of
\(\nu_0=p^{e_{\mathfrak p}-2}t\) is
\[
 t=\varepsilon(a'_{\mathfrak f}/p^{e_{\mathfrak p}})
                  c'(n')^3d^{-2}(f/p)^{-2}.
\]
At \(\mathfrak q\mid\mathfrak c'\), put
\(k=v_{\mathfrak q}(\mathfrak n')\).  The spherical valuation is \(1+3k\),
and its unit exponent is \(2\) modulo \(6\).  The cube unit from
\((n'/q^k)^3\) disappears from the symbol.  A prime of \(\mathfrak n'\)
outside \(\mathfrak c'\) has unit exponent divisible by \(6\).
Consequently the exact grouped exterior factor on
\((\mathfrak d,\mathfrak h\mathfrak n')=
  (\mathfrak c',\mathfrak f)=(\mathfrak d,\mathfrak c')=1\) is
\begin{equation}\label{eq:reflection-exterior-factorization}
\begin{split}
 \mathcal E_{\mathfrak t}(\mathfrak d,\mathfrak c')
 ={}&\frac{I_{\mathfrak h,\mathfrak f}}
          {q_{\mathfrak c'}^{1/2}q_{\mathfrak n'}}\\
 &{}\times\prod_{\mathfrak p\mid\mathfrak d}
       \omega_{\mathfrak p}
       \chi_{\mathfrak p}\!\left(
          \varepsilon a'_{\mathfrak f}c'(n')^3(d/p)^{-2}f^{-2}\right)\\
 &{}\times\prod_{\mathfrak p\mid\mathfrak f}
       s_{\mathfrak p}
       \chi_{\mathfrak p}\!\left(
          \varepsilon(a'_{\mathfrak f}/p^{e_{\mathfrak p}})
          c'(n')^3d^{-2}(f/p)^{-2}\right)^{b_{\mathfrak p}}\\
 &{}\times\prod_{\mathfrak q\mid\mathfrak c'}
       g_{2,\mathfrak q}
       \chi_{\mathfrak q}\!\left(
          \varepsilon a'_{\mathfrak f}(c'/q)(df)^{-2}\right)^2 .
\end{split}
\end{equation}
The factor \(I_{\mathfrak h,\mathfrak f}\) is present whether or not an
inactive prime occurs in \(\mathfrak c'\) or \(\mathfrak n'\).  Thus the
formula includes their overlap at such a prime as well as every other
\(\mathfrak c'\)--\(\mathfrak n'\) overlap.  The local zeros
\(s_{\mathfrak p}=0\) keep the broader active numerator support exact.

We identify the only factor involving both row and column variables.
At a pair \(\mathfrak p\mid\mathfrak d\),
\(\mathfrak q\mid\mathfrak c'\), the row factor supplies
\(\chi_{\mathfrak p}(q)\), and the denominator \(d^{-2}\) in the
spherical column factor supplies \(\chi_{\mathfrak q}(p)^{-4}\).
Allowed-generator reciprocity gives
\begin{equation}\label{eq:reflection-one-minus-four}
 \chi_{\mathfrak p}(q)\chi_{\mathfrak q}(p)^{-4}
      =\chi_{\mathfrak p}(q)^{1-4}
      =\chi_{\mathfrak p}(q)^3.
\end{equation}
Multiplying over the pairs gives \(\chi_{\mathfrak d}(c')^3\).
There is no further pairing from a prime occurring only in
\(\mathfrak n'\), since its unit exponent is divisible by \(6\).
The separated \(S\)-character is
\[
 L(\nu_0)^{-1}
  =L(a'_{\mathfrak f})^{-1}L(c')^{-1}L(n')^{-3}L(d)^2L(f)^2,
\]
because \(L(\varepsilon)=1\).  Comparing
\Cref{eq:reflection-exterior-factorization} with
\Cref{eq:reflection-frozen-factor,eq:reflection-row-factor,eq:reflection-column-factor}
therefore gives, on the indicated class slices,
\begin{equation}\label{eq:reflection-factor-separation}
 \mathcal E_{\mathfrak t}(\mathfrak d,\mathfrak c')L(\nu_0)^{-1}
  =\frac{\omega_{\upsilon,\mathfrak t}
          \rho_{\upsilon,\mathfrak t,s_d}(\mathfrak d)
          \lambda_{\upsilon,\mathfrak t,s_c}(\mathfrak c')}
         {q_{\mathfrak c'}^{1/2}q_{\mathfrak n'}}
       \chi_{\mathfrak d}(c')^3.
\end{equation}
The local bounds prove the three coefficient bounds in the lemma.
The zero extension of the quadratic symbol makes the right side zero
when \(\mathfrak d\) and \(\mathfrak c'\) meet, as the valuation-one
support requires.  A prime shared by \(\mathfrak d\) and \(\mathfrak n'\)
must instead be removed by the separate row support of \(\rho\).
Together with \((\mathfrak n',\mathfrak f)=1\) and the column support of
\(\lambda\), these conditions reproduce exactly
\((\mathfrak c'\mathfrak n',\mathfrak d\mathfrak f)=1\).

It remains to evaluate the \(S\)-integral.  The grouped digit sums
\(\mathcal R_{\mathfrak p,j,e}(t)\) are unchanged under \(t\mapsto\gamma t\)
for \(\gamma\in\Gamma\), by \Cref{eq:reflection-local-monomial}.
The same is true of each ordinary spherical factor and of the zero-digit
factor, since \(\chi_{\mathfrak p}(\gamma)=1\).  This includes active
\(j=0\) terms and the zero digit.  Thus the grouped exterior coefficient
may be taken outside the signed \(\Gamma\)-orbit.  That orbit is absolutely
dominated by \Cref{eq:reflection-reflected-absolute}.

The reflected \(S\)-factor is
\(\mathcal C_S(\gamma\nu_0,y^{-1})f_{w,\upsilon}(\gamma\nu_0/y)\).
In the legitimate substitution \(v=\gamma\nu_0/y\),
\([y]=[\nu_0][v]^{-1}\) and
\(L(y)^{-1}=L(\nu_0)^{-1}L(v)\).  Alternation and
\Cref{eq:reflection-S-cocycle} give
\[
 \mathcal C_S(\gamma\nu_0,y^{-1})
   =H_S(\gamma\nu_0,v/(\gamma\nu_0))^2
   =H_S(\nu_0,v)^2.
\]
We obtain the exact identity
\begin{equation}\label{eq:reflection-S-unfolding}
\begin{split}
 &\sum_{\gamma\in\Gamma}\int_{\Omega_l}
   \one_B([y])L(y)^{-1}\mathcal C_S(\gamma\nu_0,y^{-1})
       f_{w,\upsilon}(\gamma\nu_0/y)\,\dd_l^\times y\\
 &\quad=L(\nu_0)^{-1}\int_{X_{|\nu_0|_S/l}}
   \one_B([\nu_0][v]^{-1})H_S(\nu_0,v)^2L(v)f_{w,\upsilon}(v)
       \,\dd_{|\nu_0|_S/l}^\times v .
\end{split}
\end{equation}
The sheet norm is
\[
 \frac{|\nu_0|_S}{l}
   =\frac{q_{\mathfrak a'_{\mathfrak f}}q_{\mathfrak c'}q_{\mathfrak n'}^3}
          {q_{\mathfrak d}^2q_{\mathfrak f}^2l};
\]
there is no norm Jacobian.  Inside the integral \(\nu_0\) occurs only
through its class in \(\mathcal V\).  On the supports of the two class
indicators this class is exactly
\(\kappa_6(\mathfrak t;s_d,s_c)\).  Thus no continuous shape of the row
or column generator remains inside the profile.

Insert the two class partitions, use
\Cref{eq:reflection-factor-separation}, and split the fixed \(S\)-function
into its \(\ell\)-terms.  Putting \(l=T/Z\) in
\Cref{eq:reflection-exact-average} produces precisely the combined profile
in \Cref{eq:reflection-profiles}, and hence
\Cref{eq:reflection-global-expansion}.  For a fixed row class and fixed
ideal data \((\mathfrak h,\mathfrak f,\mathfrak a'_{\mathfrak f},\mathfrak n')\),
there are only \(|U_6||B|^2\) choices of the frequency unit and two slices.
After those choices, the class
\(\kappa_6\) is determined and there is one combined profile.  The active
sets and the ideal sums are the explicit outer sums in the expansion.

We finish by verifying the uniform absolute bound after grouping.  For
\(\sigma\ge0\),
\begin{equation}\label{eq:reflection-frozen-Dirichlet}
 E_{\mathfrak p}(\sigma)
 :=\sum_{e\ge0}\beta_{\mathfrak p}(e)q_{\mathfrak p}^{-\sigma e}
 =1+\frac{q_{\mathfrak p}^{-\sigma}
             +q_{\mathfrak p}^{-1/2-3\sigma}}
            {1-q_{\mathfrak p}^{-1-3\sigma}}
 \le 3+\sqrt2.
\end{equation}
In particular, for every \(\delta_1>0\),
\(\prod_{\mathfrak p\mid\mathfrak f}E_{\mathfrak p}(\sigma)
 \ll_{F,\delta_1}q_{\mathfrak f}^{\delta_1}\), uniformly for
\(\sigma\ge0\).  Absorb the finitely many primes of norm below
\((3+\sqrt2)^{1/\delta_1}\) into the constant and bound each remaining
factor by \(q_{\mathfrak p}^{\delta_1}\).

Put \(Q=q_{\mathfrak a'_{\mathfrak f}}q_{\mathfrak c'}q_{\mathfrak n'}^3\).
Dropping squarefreeness and the unfrozen coprimality restrictions in the
nonnegative majorant gives
\begin{equation}\label{eq:reflection-reflected-Dirichlet}
\begin{split}
 \mathcal D_{\mathfrak f}(\sigma)
 &:=\sum_{\mathfrak a'_{\mathfrak f},\mathfrak c',\mathfrak n'}
       \beta_{\mathfrak f}(\mathfrak a'_{\mathfrak f})
       q_{\mathfrak c'}^{-1/2}q_{\mathfrak n'}^{-1}Q^{-\sigma}\\
 &\le\zeta_F^S(\sigma+1/2)\zeta_F^S(3\sigma+1)
       \prod_{\mathfrak p\mid\mathfrak f}E_{\mathfrak p}(\sigma)
 <\infty\qquad(\sigma>1/2).
\end{split}
\end{equation}
The frozen ideal still ranges over \(\mathcal A(\mathfrak f)\).
For \(t_{\mathfrak d}=q_{\mathfrak d}^2q_{\mathfrak f}^2/Z\), the
absolute profile above dominates every class multiplier and yields
\begin{equation}\label{eq:reflection-reflected-Fubini}
\begin{split}
 &\int_0^\infty|\mathcal K(T)|
   \sum_{\mathfrak a'_{\mathfrak f},\mathfrak c',\mathfrak n'}
     \beta_{\mathfrak f}(\mathfrak a'_{\mathfrak f})
     q_{\mathfrak c'}^{-1/2}q_{\mathfrak n'}^{-1}
     H^{\mathrm{abs}}_{w,\upsilon}\!\left(
        \frac{Q}{t_{\mathfrak d}T}\right)\frac{\dd T}{T}\\
 &\hspace{2cm}\ll_\sigma
       t_{\mathfrak d}^{\sigma}M_\sigma(\mathcal K)
             \mathcal D_{\mathfrak f}(\sigma)<\infty.
\end{split}
\end{equation}
This justifies all exchanges in the grouped expansion.  Its constants
are uniform in the ideal data, apart from the displayed product of local
majorants.  Since \(\mathcal F(\mathfrak h)\) is finite for each fixed row,
the full expansion is absolutely convergent.
\end{proof}

The case \(m=0\) remains the separate identity
\(B_0(Z)=P_G(1/Z)\) from the coefficient-family section.

The expansion also displays the length used in the second moment.  On dyads
\(P\le q_{\mathfrak h}<2P\) and \(D\le q_{\mathfrak d}<2D\), put
\begin{equation}\label{eq:reflection-column-cutoff}
\begin{aligned}
 Q&=q_{\mathfrak a'_{\mathfrak f}}q_{\mathfrak c'}q_{\mathfrak n'}^3,
 &\qquad
 t_{\mathfrak d}&=\frac{q_{\mathfrak d}^2q_{\mathfrak f}^2}{Z},\\
 \mathcal L&=\frac{8D^2P}{Z},
 & R_0&=Z^\delta\mathcal L,
\end{aligned}
\end{equation}
where \(\delta>0\) will be chosen small.  By
\Cref{eq:reflection-active-norm}, \(t_{\mathfrak d}\le\mathcal L\).
We shall retain \(Q\le R_0\), which in particular gives
\(q_{\mathfrak c'}\le R_0\).  If \(R_0<1\), this range is empty because
\(Q\) is the norm of an integral ideal.  The tail estimate below justifies
this truncation for every dyad, including that case.

\subsection{The quadratic large sieve and the moment estimate}

We use the quadratic large sieve established over \(\Q\) by Heath-Brown
and extended to Hecke families by Goldmakher--Louvel
\cite{HeathBrown95,GoldmakherLouvel13}.  We state the latter result in the
form needed for \Cref{eq:reflection-global-expansion} and verify its family
hypotheses for the present generator convention.

\begin{lemma}[Quadratic large sieve for the exterior symbols]
\label{lem:reflection-quadratic-sieve}
For $C,D\geq1$ and complex coefficients $z_{\mathfrak c}$ supported on
squarefree exterior ideals with $q_{\mathfrak c}\leq C$, one has, for every
$\epsilon>0$,
\begin{equation}\label{eq:reflection-quadratic-sieve}
 \sum_{\substack{\mathfrak d\ \mathrm{squarefree}\\q_{\mathfrak d}\leq D}}
   \left|\sum_{\substack{\mathfrak c\ \mathrm{squarefree}\\q_{\mathfrak c}\leq C}}
       z_{\mathfrak c}\chi_{\mathfrak d}(c)^3\right|^2
 \ll_{F,S,\epsilon}(C+D)(CD)^\epsilon
             \sum_{\mathfrak c}|z_{\mathfrak c}|^2.
\end{equation}
Arbitrary restrictions on the rows and separate restrictions or unitary
factors on the columns are allowed by deleting rows and changing the
coefficients.
\end{lemma}

\begin{proof}
We verify Definition~1 of Goldmakher--Louvel before applying their
Theorem~1.1 \cite{GoldmakherLouvel13}.  Choose an ordinary integral ideal
$\mathfrak q_S$ supported on the finite places of $S$, with a positive
exponent $m_v$ at each such place large enough that
$1+\mathfrak p_v^{m_v}\mathcal O_v\subset F_v^{\times2}$.
Every local quadratic character at $v$ is trivial on this subgroup, so its
conductor exponent is at most $m_v$.  Exterior ideals are exactly the
ordinary ideals prime to $\mathfrak q_S$, with the same norms.

For a squarefree exterior ideal $\mathfrak c$, let $\Xi_{\mathfrak c}$ be
the primitive Hecke character of the possibly trivial quadratic Kummer extension
$F(\sqrt c)/F$.  It is independent of the allowed generator, since two
such generators differ by a sixth power of an $S$-unit.  At an exterior
prime its local conductor exponent is one if the prime divides
$\mathfrak c$ and zero otherwise: at the odd residue characteristics under
consideration an odd valuation gives a tamely ramified quadratic extension,
whereas a unit gives an unramified or trivial extension.  Thus
\[
 \operatorname{cond}(\Xi_{\mathfrak c})
       =\mathfrak c\mathfrak s_{\mathfrak c},
 \qquad \mathfrak s_{\mathfrak c}\mid\mathfrak q_S.
\]
Its infinite type is trivial because all infinite places of $F$ are complex.
On an exterior ideal $\mathfrak d$,
\[
 \Xi_{\mathfrak c}(\mathfrak d)=\chi_{\mathfrak d}(c)^3.
\]
This includes the zero at a shared prime, which is in the conductor of
$\Xi_{\mathfrak c}$.  These are the primitive global characters, not a
character with additional artificially deleted Euler factors.

Let $B_2$ be the image of $B$ in $F_S^\times/F_S^{\times2}$ and set
$\mathrm{cl}_2(\mathfrak c)=[c]_{S,2}\in B_2$.  Allowed generators
multiply modulo $U_S^6$, so this extends to a homomorphism on the group of
all fractional exterior ideals.  It has finite image.  It is also ray-class
data for $\mathfrak q_S$: for a principal exterior ideal $(x)$ its value is
the $B_2$-projection of $[x]_{S,2}$ in the direct sum induced by $E\oplus B$;
if $x\equiv1\pmod{\mathfrak q_S}$, all its local classes at $S$ are squares,
so this projection is trivial.

For coprime squarefree exterior ideals, the exact reciprocity of allowed
generators gives
\[
 \Xi_{\mathfrak c}(\mathfrak d)
     =\chi_{\mathfrak d}(c)^3
     =\chi_{\mathfrak c}(d)^3
     =\Xi_{\mathfrak d}(\mathfrak c).
\]
Thus the finite-class reciprocity factor in that definition may be taken
identically one.  If $\mathfrak c_1,\mathfrak c_2$ are coprime squarefree exterior ideals and have the
same $\mathrm{cl}_2$, their quotient $c_1/c_2$ is a square at every place
of $S$.  The primitive character inducing
$\Xi_{\mathfrak c_1}\overline{\Xi_{\mathfrak c_2}}$ is consequently
trivial there.  At every prime of $\mathfrak c_1\mathfrak c_2$ it is tamely
ramified with conductor exponent one, and at every other exterior prime it
is unramified.  Its \emph{global} conductor is therefore exactly
$\mathfrak c_1\mathfrak c_2$.  The product here is understood as the
primitive inducing character, with the convention specified in
\cite[\S2]{GoldmakherLouvel13}.  This verifies the conductor-product axiom,
including the part at $S$.

We have thus obtained a quadratic Hecke family in the precise sense of
\cite[Definition~1]{GoldmakherLouvel13}.  Its Theorem~1.1 gives
\Cref{eq:reflection-quadratic-sieve}.  Apply that theorem to the full
family and then delete unwanted rows or set unwanted coefficients to zero;
a single class slice need not itself be a Hecke family.
\end{proof}

\begin{proof}[Proof of \Cref{prop:reflection}]
The initial element-to-ideal reduction lets us fix one row-unit class
\(\upsilon\in U_6\) and sum over ideal rows at a bounded cost.
Fix dyads $P\leq q_{\mathfrak h}<2P$ and
$D\leq q_{\mathfrak d}<2D$ which meet the row set.  They satisfy
\begin{equation}\label{eq:reflection-PD-range}
 PD\ll_{\mathscr K}Z^{M_0},\qquad P,D\geq1.
\end{equation}
There are $O_F(P^{1/2})$ powerful ideals on the first dyad.  In fact every
powerful ideal has a unique expression $\mathfrak r^2\mathfrak s^3$ with
$\mathfrak s$ squarefree (choose its exponent as the parity of each prime
exponent).  The fixed-field bound for the number of ideals of norm at most
$t$ is $O_F(t)$, hence
\begin{equation}\label{eq:reflection-powerful-count}
 \#\{\mathfrak h\text{ powerful}:q_{\mathfrak h}<2P\}
 \ll_F P^{1/2}\sum_{\mathfrak s}q_{\mathfrak s}^{-3/2}
 \ll_F P^{1/2}.
\end{equation}

Fix \(\mathfrak h\), and use \(Q,t_{\mathfrak d},\mathcal L,R_0\) from
\Cref{eq:reflection-column-cutoff}.  For each
\(\mathfrak f\in\mathcal F(\mathfrak h)\) and
\(\mathbf s=(u,s_d,s_c)\in U_6\times B^2\), split its partial sum in
\Cref{eq:reflection-global-expansion} according to \(Q\le R_0\) or
\(Q>R_0\).  Denote these sums, which still include
\(\mathfrak a'_{\mathfrak f},\mathfrak n',\mathfrak c'\), by
\(\operatorname{Ret}_{\mathfrak f,\mathbf s,\mathfrak d}\) and
\(\operatorname{Tail}_{\mathfrak f,\mathbf s,\mathfrak d}\).
Sum them over the admitted \(\mathfrak f\) and the finite \(\mathbf s\)
to define \(\operatorname{Ret}_{\mathfrak h,\mathfrak d}\) and
\(\operatorname{Tail}_{\mathfrak h,\mathfrak d}\).  Thus, on a row with
class \(\upsilon\),
\[
 B_m(Z)=\operatorname{Ret}_{\mathfrak h,\mathfrak d}
          +\operatorname{Tail}_{\mathfrak h,\mathfrak d}.
\]
Set both parts to zero outside the squarefree rows
\(\mathfrak d\) coprime to \(\mathfrak h\).  The threshold \(R_0\) is
common to the dyad and independent of the individual row.  When \(R_0<1\)
the retained part is empty, so no large sieve at a length below one is
required.

We first bound the discarded terms, uniformly in this case as well.  Put
$\mathcal J=\max(Z^\delta,\mathcal L^{-1})$, and fix $N>\sigma>1/2$.
For $Q>R_0$ we have both $Q/t_{\mathfrak d}>Z^\delta$ and
$Q/t_{\mathfrak d}\geq\mathcal L^{-1}$, because $Q\geq1$.
The profile estimate therefore gives
\[
 \one_{Q>R_0}|\Phi_{\upsilon,\kappa}(Q/t_{\mathfrak d})|
 \ll_N \mathcal J^{-(N-\sigma)}(t_{\mathfrak d}/Q)^\sigma.
\]
Summing absolutely with \Cref{eq:reflection-reflected-Dirichlet}, for one
active product and one fixed \(\mathbf s\), yields
\begin{equation}\label{eq:reflection-tail-bound}
 |\operatorname{Tail}_{\mathfrak f,\mathbf s,\mathfrak d}|
 \ll_{N,\sigma,F}
 \mathcal L^\sigma\mathcal J^{-(N-\sigma)}
 \zeta_F^S(\sigma+1/2)\zeta_F^S(3\sigma+1)
 \prod_{\mathfrak p\mid\mathfrak f}E_{\mathfrak p}(\sigma).
\end{equation}
This one estimate includes all column, cube, and frozen-exponent tails.
If $R_0<1$, then $\mathcal J=\mathcal L^{-1}$ and the corresponding
entire partial sum is $O(\mathcal L^N\prod E_{\mathfrak p}(\sigma))$.

For every $\delta_1>0$, \Cref{eq:reflection-frozen-Dirichlet} also gives
\begin{equation}\label{eq:reflection-active-choice-sum}
 \sum_{\mathfrak f\mid\rad(\mathfrak h)}
       \prod_{\mathfrak p\mid\mathfrak f}E_{\mathfrak p}(\sigma)
   =\prod_{\mathfrak p\mid\mathfrak h}(1+E_{\mathfrak p}(\sigma))
   \ll_{F,\delta_1}q_{\mathfrak h}^{\delta_1}
   \qquad(\sigma\geq0).
\end{equation}
The same finite-small-prime argument proves this inequality; including all
divisors only adds active choices that may actually be forbidden.  Sum
\Cref{eq:reflection-tail-bound} over them and over the fixed finite class
list, square, and use $O_F(D)$ for the number of row ideals on the second
dyad and \Cref{eq:reflection-powerful-count}.  The total squared tail is
\begin{equation}\label{eq:reflection-tail-mass}
 \ll_{N,\sigma,F,\delta_1}
 P^{1/2+2\delta_1}D\,
 \mathcal L^{2\sigma}\mathcal J^{-2(N-\sigma)}.
\end{equation}
Here \Cref{eq:reflection-PD-range} implies
$P\ll Z^{M_0}$ and $\mathcal L\ll Z^{2M_0-1}$, while
$\mathcal J\geq Z^\delta$.  Given any $A_0>0$, choosing $N$ with
\[
 2\delta(N-\sigma)>A_0+2M_0\delta_1+2\sigma(2M_0-1)
\]
makes \Cref{eq:reflection-tail-mass}
$O(Z^{-A_0}P^{1/2}D)$.  This proves the promised absolute domination of
the discarded range, including all dyads for which $R_0<1$.

It remains to bound the retained terms when \(R_0\ge1\).  Fix
\(\mathfrak f,\mathbf s,\mathfrak a'_{\mathfrak f},\mathfrak n'\).
Then \(\kappa=\kappa_6(\mathfrak t;s_d,s_c)\) is fixed, while
\(\rho_{\upsilon,\mathfrak t,s_d}\) and
\(\lambda_{\upsilon,\mathfrak t,s_c}\) are coherent functions of all rows
and columns in the two slices.  Split \(\mathfrak c'\) into dyads
\(C\le q_{\mathfrak c'}<2C\), \(C\ge1\).  Put
\[
 v_0=\frac{Zq_{\mathfrak a'_{\mathfrak f}}q_{\mathfrak n'}^3C}
                {D^2q_{\mathfrak f}^2},\qquad
 x=\log(q_{\mathfrak c'}/C)-2\log(q_{\mathfrak d}/D).
\]
Then $x\in[-2\log2,\log2]$, and the profile in
\Cref{eq:reflection-global-expansion} is $\Phi_{\upsilon,\kappa}(v_0e^x)$.
Choose a fixed $\varphi\in C_c^\infty(\R)$ equal to one on that interval
and set $g_{v_0}(x)=\varphi(x)\Phi_{\upsilon,\kappa}(v_0e^x)$.  Applying
\Cref{eq:sheet-fourier-separation} with this cutoff gives
\begin{equation}\label{eq:reflection-Fourier-separation}
 \int_\R|\widehat g_{v_0}(\tau)|\,\dd\tau
 \ll_N\min(v_0^\alpha,v_0^{-N})\ll_N1.
\end{equation}
Fourier inversion now separates the norm factors as
$(q_{\mathfrak c'}/C)^{i\tau}(q_{\mathfrak d}/D)^{-2i\tau}$ with the
bounded integral cost in \Cref{eq:reflection-Fourier-separation}.
The retained indicator
$\one_{q_{\mathfrak a'_{\mathfrak f}}q_{\mathfrak c'}q_{\mathfrak n'}^3\leq R_0}$
is a column coefficient after the two frozen ideals are fixed.

Use Minkowski's inequality for this Fourier integral and apply
\Cref{lem:reflection-quadratic-sieve} to each $\tau$.  Separate row factors
have modulus at most one, and the column square mass is
\[
 \sum_{C\leq q_{\mathfrak c'}<2C}q_{\mathfrak c'}^{-1}\ll_F1
\]
by ideal counting.  For any $\delta_2>0$, the row $\ell^2$ norm of this
one column dyad, before the factor
$\beta_{\mathfrak f}(\mathfrak a'_{\mathfrak f})/q_{\mathfrak n'}$, is
therefore at most
\begin{equation}\label{eq:reflection-one-column-dyad}
 \ll_{N,\delta_2}(D+C)^{1/2}(DC)^{\delta_2/2}
             \min(v_0^\alpha,v_0^{-N}).
\end{equation}
The restrictions $(\mathfrak d,\mathfrak h\mathfrak n')=1$ merely delete
rows.  The zero extension of the symbol already treats shared row and
column primes, so no joint restriction was removed in applying the sieve.

For the retained range, $C\leq R_0$ and
$q_{\mathfrak n'}\leq R_0^{1/3}$.  There are $O(\log(2R_0))$ such column
dyads, and partial summation of the fixed-field ideal count gives
\[
 \sum_{q_{\mathfrak n'}\leq R_0^{1/3}}q_{\mathfrak n'}^{-1}
       \ll_F\log(2R_0).
\]
Drop the final minimum in \Cref{eq:reflection-one-column-dyad} and use
Minkowski successively for the column dyads, $\mathfrak n'$, and
$\mathfrak a'_{\mathfrak f}$.
\Cref{eq:reflection-frozen-Dirichlet} at $\sigma=0$ shows that their
retained row norm, for fixed \(\mathfrak h,\mathfrak f,\mathbf s\), is
\[
 \ll_{\delta_2}
 \left(\prod_{\mathfrak p\mid\mathfrak f}E_{\mathfrak p}(0)\right)
 \log^2(2R_0)(D+R_0)^{1/2}(DR_0)^{\delta_2/2}.
\]
Sum these norms over the admitted \(\mathfrak f\) and the finite
\(\mathbf s\), using \Cref{eq:reflection-active-choice-sum} for the former.
Square the resulting estimate for each fixed \(\mathfrak h\), and only
then sum the squares over \(\mathfrak h\).  Since each row has a unique powerful part, no triangle
inequality across different powerful parts is necessary.
\Cref{eq:reflection-powerful-count} yields
\begin{equation}\label{eq:reflection-retained-mass}
 \sum_{\substack{\mathfrak h\ \mathrm{powerful}\\P\leq q_{\mathfrak h}<2P}}
 \sum_{D\leq q_{\mathfrak d}<2D}
   |\operatorname{Ret}_{\mathfrak h,\mathfrak d}|^2
 \ll_{\delta_1,\delta_2}
 P^{1/2+2\delta_1}\log^4(2R_0)(D+R_0)(DR_0)^{\delta_2}.
\end{equation}
All the finite unit and $S$-class choices have only multiplied this by a
fixed constant.

We make the power losses explicit.  On the range
\Cref{eq:reflection-PD-range},
\[
 D+R_0\ll Z^\delta\left(D+\frac{D^2P}{Z}\right),\qquad
 DR_0\ll Z^{3M_0-1+\delta},\qquad P\ll Z^{M_0}.
\]
For any $\delta_3>0$, $\log^4(2R_0)\ll_{\delta_3}Z^{\delta_3}$ because
$1\leq R_0\ll Z^{2M_0-1+\delta}$.  Choose the positive auxiliaries so that
\[
 \delta+2M_0\delta_1+(3M_0-1+\delta)\delta_2+\delta_3<\epsilon.
\]
Then \Cref{eq:reflection-retained-mass} is bounded by the right side of
\Cref{eq:reflection-dyadic-mass}.  The retained mass is zero for $R_0<1$,
and \Cref{eq:reflection-tail-mass} is smaller than that right side in all
cases after $N$ is chosen as above.  Adding the two parts and restoring the
bounded element multiplicities proves \Cref{eq:reflection-dyadic-mass}.

Finally, for every nonempty dyad, $P\geq1$ and
\Cref{eq:reflection-PD-range} give
\[
 P^{1/2}D\leq PD\ll Z^{M_0}\leq Z^{2M_0-1},\qquad
 P^{1/2}\frac{D^2P}{Z}
       =\frac{(PD)^2}{ZP^{1/2}}\ll Z^{2M_0-1},
\]
where the first comparison uses $M_0>1$.  There are only $O((\log Z)^2)$
possible dyads.  Applying the dyadic estimate with a smaller preliminary
loss absorbs this factor and proves \Cref{eq:reflection-total-mass}.
Multiplying $\mathscr K$ by any element of a fixed compact norm-one set
only enlarges the fixed compact shape set used in the initial multiplicity
bound, which proves the asserted uniformity.
\end{proof}
